
\section{AA-02: Zeeman slowing and the Doppler limit}
\source{2025 PHYS4055 Exam, Question 2(d)--(e);
Week 10 tutorial \texttt{Tutorial\_questions.pdf}, Question 2.}
\subsection*{Given and modelling choice}
Consider the 2025 Yb beam:
\(M=1.73\times10^{-25}\) kg, source temperature
\(T=300\) K, cooling wavelength \(\lambda=399\) nm,
and excited-state lifetime \(\tau=5.5\) ns.
The question does not define a unique initial velocity,
so explicitly choose the 3D rms thermal speed
\(v_{\rm rms}=\sqrt{3k_BT/M}\).
This gives
\[
v_0=\sqrt{\frac{3(1.38065\times10^{-23})(300)}
{1.73\times10^{-25}}}
=\boxed{268\ {\rm m\,s^{-1}}}.
\]
The most probable speed \(\sqrt{2k_BT/M}\) would give a
different stopping distance; state which convention you use.

\subsection*{Step 1: maximum photon-scattering force}
A saturated cycling transition can scatter at most
approximately \(\Gamma/2\) photons per second,
\(\Gamma=1/\tau\). One absorbed photon transfers
momentum \(p_\gamma=h/\lambda\) opposite the atomic beam:
\[
F_{\rm max}=\frac{h}{2\lambda\tau},\quad
a_{\rm max}=\frac{h}{2M\lambda\tau}.
\]
The question specifies \emph{half} the maximum
deceleration:
\[
\boxed{a=\frac{a_{\rm max}}2
=\frac{h}{4M\lambda\tau}
\simeq4.36\times10^{5}\ {\rm m\,s^{-2}}}.
\]

\subsection*{Step 2: stopping length and photon count}
Using \(v_f^2=v_0^2-2aL\) with \(v_f=0\):
\[
\boxed{L=\frac{v_0^2}{2a}
\simeq0.0823\ {\rm m}=8.23\ {\rm cm}}.
\]
Each counter-propagating photon changes the
atom's momentum by \(h/\lambda\), so the required
net momentum \(Mv_0\) corresponds to
\[
\boxed{N_{\gamma}\simeq\frac{Mv_0}{h/\lambda}
=\frac{Mv_0\lambda}{h}
\simeq2.79\times10^4\ \text{photons}}.
\]
Spontaneous emission has zero \emph{mean} momentum
along the slowing axis but causes recoil heating.

\subsection*{Step 3: Doppler cooling and Zeeman resonance}
The two-level Doppler cooling limit is
\[
\boxed{T_D=\frac{\hbar\Gamma}{2k_B}
=\frac{h}{4\pi k_B\tau}
\simeq6.94\times10^{-4}\ {\rm K}=0.694\ {\rm mK}}.
\]
In optical molasses, red-detuned counter-propagating
beams are seen closer to resonance by an atom moving
against them, creating velocity-dependent damping.
As an atom slows, its Doppler shift changes; a
Zeeman slower uses a position-dependent magnetic field
to shift the atomic transition and maintain resonance.
For approximately constant deceleration,
\(v(z)=\sqrt{v_0^2-2az}\), motivating a
square-root magnetic-field profile when the Zeeman
shift is roughly linear in \(B\).

\checkmethod{Scattering force, slowing distance and Doppler
temperature depend on different quantities. A small
lifetime increases maximum force but also increases
the simple two-level Doppler limit. Include the factor
of 1/2 when the question asks for half maximum acceleration.}
